\section{\texorpdfstring{هفته‌ی 4: برآورد بیشینه درست‌نمایی \lr{(Maximum
Likelihood)}}{هفته‌ی 4: برآورد بیشینه درست‌نمایی (Maximum Likelihood)}}

\begin{quote}
\textbf{فصل 2: روش‌های برآورد پارامتر} · درس 1 ساعت · \lr{lab} 2 ساعت
\textbf{داده‌ی اصلی:} آزمایش فراموشی هفته‌های 2\lr{--}3 (شمارش‌های یادآوری
دوجمله‌ای)، به‌علاوه‌ی یک آزمایش \lr{SDT} جدید با سه شرط پاداش
\lr{(payoff).}
\end{quote}

\begin{center}\rule{0.5\linewidth}{0.5pt}\end{center}

\subsection{واژه‌های کلیدی}

واژه‌های مهم این درس با معادل فارسی‌شان (حتی اگر در درس‌های قبل آمده
باشند).

\begin{longtable}[]{@{}
  >{\raggedright\arraybackslash}p{(\linewidth - 4\tabcolsep) * \real{0.3333}}
  >{\raggedright\arraybackslash}p{(\linewidth - 4\tabcolsep) * \real{0.3333}}
  >{\raggedright\arraybackslash}p{(\linewidth - 4\tabcolsep) * \real{0.3333}}@{}}
\caption{\centering واژه‌های کلیدی}\tabularnewline
\toprule\noalign{}
\begin{minipage}[b]{\linewidth}\raggedright
\lr{English}
\end{minipage} & \begin{minipage}[b]{\linewidth}\raggedright
فارسی
\end{minipage} & \begin{minipage}[b]{\linewidth}\raggedright
توضیح کوتاه
\end{minipage} \\
\midrule\noalign{}
\endfirsthead
\toprule\noalign{}
\begin{minipage}[b]{\linewidth}\raggedright
\lr{English}
\end{minipage} & \begin{minipage}[b]{\linewidth}\raggedright
فارسی
\end{minipage} & \begin{minipage}[b]{\linewidth}\raggedright
توضیح کوتاه
\end{minipage} \\
\midrule\noalign{}
\endhead
\bottomrule\noalign{}
\endlastfoot
\lr{likelihood} & درست‌نمایی & احتمال داده به‌عنوان تابعی از پارامترها \\
\lr{estimate} / \lr{estimator} & برآورد / برآوردگر & مقدار به‌دست‌آمده
برای پارامتر / روش به دست آوردن آن \\
\lr{parameter} & پارامتر & کمیت قابل تنظیم مدل که بین افراد یا شرایط فرق
می‌کند \\
\lr{forgetting curve} & منحنی فراموشی & کاهش یادآوری با گذشت زمان \\
\lr{binomial distribution} & توزیع دوجمله‌ای & توزیع تعداد موفقیت در
\lr{n} کوشش \\
\lr{signal detection theory (SDT)} & نظریه‌ی تشخیص علامت & مدلی که حساسیت
را از سوگیری پاسخ جدا می‌کند \\
\lr{log-likelihood} / \lr{NLL} & لگاریتم درست‌نمایی / منفی آن & مقداری که
بهینه‌سازها کمینه یا بیشینه می‌کنند \\
\lr{least squares} & کمترین مربعات & روش برازش با کمینه کردن مجموع مربع
خطاها \\
\lr{maximum likelihood estimate (MLE)} & برآورد بیشینه‌ی درست‌نمایی &
پارامترهایی که داده را محتمل‌ترین می‌کنند \\
\lr{noise} & نویز & تغییرات تصادفی در شواهد یا در داده \\
\lr{fit} / \lr{fitting} & برازش & تنظیم پارامترها تا مدل با داده جور
شود \\
\lr{deviance (G²)} & انحراف \lr{(deviance)} & معیار عدم برازش نسبت به
مدل اشباع‌شده \\
\lr{bounds} / \lr{constraints} & کران‌ها / قیدها & محدودیت‌های مجاز برای
مقدار پارامترها \\
\lr{response bias} & سوگیری پاسخ & تمایل به یک پاسخ، مستقل از محرک \\
\lr{sensitivity (d′)} & حساسیت & فاصله‌ی توزیع علامت و نویز \\
\lr{trial} & کوشش \lr{(trial)} & یک بار انجام تکلیف توسط شرکت‌کننده \\
\lr{optimizer} & بهینه‌ساز & الگوریتمی که کمینه‌ی تابع را پیدا می‌کند \\
\lr{SSE} / \lr{RMSD} & مجموع مربعات خطا / ریشه‌ی میانگین مربع انحراف &
معیارهای فاصله‌ی داده و پیش‌بینی \\
\lr{saturated model} & مدل اشباع‌شده & مدلی که داده را دقیقاً بازتولید
می‌کند \\
\lr{likelihood-ratio test (LRT)} & آزمون نسبت درست‌نمایی & مقایسه‌ی دو مدل
تودرتو \\
\lr{exponential} / \lr{power function} & تابع نمایی / توانی & دو شکل
رایج برای منحنی فراموشی و یادگیری \\
\lr{liberal} / \lr{conservative} & آسان‌گیر / محافظه‌کار & تمایل به گفتن
«بله» / «نه» \\
\lr{token} & \lr{token} (واحد متن) & کلمه یا بخشی از کلمه که مدل زبانی
پیش‌بینی می‌کند \\
\lr{perplexity} & \lr{perplexity} (سردرگمی) & \lr{e} به توان \lr{NLL} هر
\lr{token} \\
\end{longtable}

\begin{center}\rule{0.5\linewidth}{0.5pt}\end{center}

\subsection{اهداف یادگیری}

در پایان این هفته، دانشجویان می‌توانند:

\begin{enumerate}
\def\labelenumi{\arabic{enumi}.}
\tightlist
\item
  احتمال (داده متغیر، پارامترها ثابت) را از درست‌نمایی (داده ثابت،
  پارامترها متغیر) \textbf{متمایز کنند}.
\item
  درست‌نمایی و لگاریتم درست‌نمایی را برای مدل‌های داده‌ی دوجمله‌ای (برنولی) و
  گاوسی \textbf{استخراج کنند} و \textbf{نشان دهند} که کمترین مربعات همان
  بیشینه درست‌نمایی تحت نویز گاوسی با واریانس برابر است.
\item
  برازش مطلق را با انحراف \lr{(deviance)} \(G^2\) \textbf{ارزیابی کنند}
  و \lr{MLE} را با برآوردهای کمترین مربعات \textbf{مقایسه کنند}.
\item
  مدل‌های مقیدشده (مثلاً «پاداش سوگیری را تغییر می‌دهد، نه حساسیت را») را
  به‌عنوان فرضیه‌های روان‌شناختی واقعی \textbf{تفسیر کنند}.
\end{enumerate}

\begin{center}\rule{0.5\linewidth}{0.5pt}\end{center}

\subsection{مقدمه}

در هفته‌ی 3، کمترین مربعات با همه‌ی نقاط داده به یک اندازه قابل اعتماد
رفتار کرد. ولی نسبت یادآوری 0.50 که روی 50 آیتم اندازه‌گیری شده
\emph{پرنویزتر} از نسبت 0.95 است: واریانس دوجمله‌ای \(f(1-f)/n\) به \(f\)
بستگی دارد. کمترین مربعات همچنین درباره‌ی داده‌هایی که اعداد پیوسته
نیستند، مثل یک پاسخ بله/نه، یک انتخاب بین چند گزینه یا یک درجه‌ی اطمینان،
حرفی برای گفتن ندارد.

بیشینه درست‌نمایی «پیش‌بینی‌ها را به داده نزدیک کن» را با سؤالی دقیق‌تر
جایگزین می‌کند:

برای جواب دادن، مدل باید بگوید \emph{هر} مجموعه‌داده‌ی ممکن چقدر محتمل
است، و این دقیقاً همان مدل نویز است (جزء 4 در کالبدشکافی هفته‌ی 2). وقتی
آن را داشته باشیم، \lr{MLE} با شمارش‌ها، انتخاب‌ها، زمان‌های واکنش و
ترکیب‌هایشان به یک شکل کار می‌کند. همچنین بدون هزینه‌ی اضافه یک معیار مطلق
عدم برازش \lr{(deviance)}، پایه‌ی خطاهای استاندارد (هفته‌ی 5)، معیارهای
اطلاعاتی (هفته‌ی 7) و جمله‌ی درست‌نمایی در قاعده‌ی بیز (هفته‌های 9\lr{--}13) را به
ما می‌دهد. \textbf{درست‌نمایی مهم‌ترین مفهوم این درس است.}

\paragraph{این موضوع کجا به کار می‌رود؟}

\begin{itemize}
\tightlist
\item
  \textbf{تقریباً هر مدل شناختی منتشرشده} با بیشینه‌ی درست‌نمایی یا نسخه‌ی
  بیزی آن (هفته‌های 9 تا 13) برازش می‌شود. وقتی مقاله‌ای «برازش مدل»،
  \lr{«AIC»} یا \lr{«log-likelihood»} گزارش می‌کند، \lr{MLE} پشت آن است.
\item
  \textbf{روان‌فیزیک \lr{(psychophysics):}} تابع روان‌سنجی (احتمال تشخیص
  یک محرک بر حسب شدت آن) در هر آزمایشگاه بینایی و شنوایی با \lr{MLE}
  دوجمله‌ای برازش می‌شود.
\item
  \textbf{آزمودن فرضیه‌ها داخل یک مدل:} «آیا پاداش حساسیت را عوض می‌کند یا
  فقط سوگیری را؟» به مقایسه‌ی دو درست‌نمایی تبدیل می‌شود (بخش 6).
\item
  \textbf{هوش مصنوعی:} مدل‌های بزرگ زبانی با بیشینه کردن درست‌نمایی متن
  آموزش می‌بینند (بخش هوش مصنوعی و \lr{LLM}ها را ببینید).
\end{itemize}

\begin{center}\rule{0.5\linewidth}{0.5pt}\end{center}

\Needspace{11\baselineskip}

\subsection{جزوه‌ی درس}

\paragraph{به زبان ساده}

تصور کنید مقدارهای مختلف پارامتر را یکی‌یکی امتحان می‌کنید. برای هر مقدار
از مدل بپرسید: «اگر این مقدار واقعی بود، داده‌ای که واقعاً جمع کرده‌ام چقدر
محتمل بود؟» این عدد \textbf{درست‌نمایی \lr{(likelihood)}} آن مقدار
پارامتر است. \textbf{برآورد بیشینه‌ی درست‌نمایی \lr{(MLE)}} مقداری از
پارامتر است که تحت آن، داده‌ی مشاهده‌شده محتمل‌ترین بود. مثال: کسی 7 کلمه
از 10 کلمه را به یاد می‌آورد. اگر احتمال واقعی یادآوری او 0.5 بود، به دست
آمدن دقیقاً 7 نسبتاً بعید بود (0.117)؛ اگر 0.7 بود، محتمل‌تر بود (0.267).
پس 0.7 توضیح بهتری است و در واقع \lr{MLE} برابر 7/10 است. در عمل با
\textbf{لگاریتم} درست‌نمایی کار می‌کنیم (تا ضرب‌ها به جمع تبدیل شوند) و از
کامپیوتر می‌خواهیم منفی آن \lr{(NLL)} را با بهینه‌سازهای هفته‌ی 3
\emph{کمینه} کند. دو جایزه هم می‌گیریم: معلوم می‌شود کمترین مربعات حالت
خاصی از \lr{MLE} است، و درست‌نمایی یک آزمون مطلق برای عدم برازش (\(G^2\))
می‌دهد.

\textbf{زمان‌بندی پیشنهادی (60 دقیقه):} احتمال در برابر درست‌نمایی 10 ·
\lr{MLE} دوجمله‌ای 15 · لگاریتم درست‌نمایی و \lr{NLL} 10 · حالت گاوسی
\lr{(LS} = \lr{MLE)} 10 · برازش مطلق \(G^2\) 10 · مدل‌های مقیدشده 5.

\subsubsection{1. احتمال در برابر درست‌نمایی}

تابع یکسان \(p(\mathbf{y}\mid\boldsymbol\theta)\) را می‌شود به دو شکل
خواند:

\begin{longtable}[]{@{}
  >{\raggedright\arraybackslash}p{(\linewidth - 8\tabcolsep) * \real{0.2000}}
  >{\raggedright\arraybackslash}p{(\linewidth - 8\tabcolsep) * \real{0.2000}}
  >{\raggedright\arraybackslash}p{(\linewidth - 8\tabcolsep) * \real{0.2000}}
  >{\raggedright\arraybackslash}p{(\linewidth - 8\tabcolsep) * \real{0.2000}}
  >{\raggedright\arraybackslash}p{(\linewidth - 8\tabcolsep) * \real{0.2000}}@{}}
\caption{\centering 1. احتمال در برابر درست‌نمایی}\tabularnewline
\toprule\noalign{}
\begin{minipage}[b]{\linewidth}\raggedright
خوانش
\end{minipage} & \begin{minipage}[b]{\linewidth}\raggedright
ثابت
\end{minipage} & \begin{minipage}[b]{\linewidth}\raggedright
متغیر
\end{minipage} & \begin{minipage}[b]{\linewidth}\raggedright
نام
\end{minipage} & \begin{minipage}[b]{\linewidth}\raggedright
جمع/انتگرالش 1 است؟
\end{minipage} \\
\midrule\noalign{}
\endfirsthead
\toprule\noalign{}
\begin{minipage}[b]{\linewidth}\raggedright
خوانش
\end{minipage} & \begin{minipage}[b]{\linewidth}\raggedright
ثابت
\end{minipage} & \begin{minipage}[b]{\linewidth}\raggedright
متغیر
\end{minipage} & \begin{minipage}[b]{\linewidth}\raggedright
نام
\end{minipage} & \begin{minipage}[b]{\linewidth}\raggedright
جمع/انتگرالش 1 است؟
\end{minipage} \\
\midrule\noalign{}
\endhead
\bottomrule\noalign{}
\endlastfoot
پیش از آزمایش & \(\boldsymbol\theta\) & \(\mathbf{y}\) & احتمال
(جرم/چگالی) داده & بله، روی \(\mathbf{y}\) \\
پس از آزمایش & \(\mathbf{y}\) (مشاهده‌شده) & \(\boldsymbol\theta\) &
\textbf{درست‌نمایی}
\(L(\boldsymbol\theta) = p(\mathbf{y}\mid\boldsymbol\theta,\mathcal{M})\)
& \textbf{نه}، روی \(\boldsymbol\theta\) نه \\
\end{longtable}

درست‌نمایی احتمال درست بودن پارامترها \textbf{نیست}. احتمال داده
\emph{تحت} هر مقدار ممکن پارامتر است. (تبدیل آن به یک توزیع احتمال روی
\(\boldsymbol\theta\) به یک پیشین نیاز دارد: بیز، هفته‌ی 9.)

\textbf{مثال.} یک شرکت‌کننده \(y = 7\) از \(n = 10\) آیتم را به یاد
می‌آورد. تحت یک مدل دوجمله‌ای با احتمال یادآوری \(\theta\):

\[
L(\theta) = \binom{10}{7}\theta^{7}(1-\theta)^{3}.
\]

\(L(0.5) = 0.117\)، \(L(0.7) = 0.267\): داده‌های مشاهده‌شده اگر
\(\theta = 0.7\) باشد حدود 2.3 برابر محتمل‌ترند تا اگر \(\theta = 0.5\)
باشد.

\paragraph{خواندن ریاضی}

\(L(\theta) = \binom{10}{7}\theta^{7}(1-\theta)^{3}\):

\begin{itemize}
\tightlist
\item
  \((1-\theta)^3\): احتمال اینکه 3 آیتم دیگر فراموش شوند.
\item
  \(\binom{10}{7} = 120\) به \(\theta\) بستگی ندارد، پس هیچ‌وقت
  \emph{محل} بیشینه را عوض نمی‌کند.
\item
  با \(\theta = 0.7\):
  \(120 \times 0.7^7 \times 0.3^3 = 120 \times 0.0824 \times 0.027 = 0.267\).
\end{itemize}

تغییر نگاه اصلی: در احتمال، \(\theta\) ثابت است و درباره‌ی داده‌های مختلف
می‌پرسیم؛ در درست‌نمایی، داده ثابت است (آن را داریم) و \(\theta\)های مختلف
را مقایسه می‌کنیم.

\subsubsection{2. برآورد بیشینه درست‌نمایی}

\[
\hat{\boldsymbol\theta}_{\text{\lr{ML}}} = \arg\max_{\boldsymbol\theta} L(\boldsymbol\theta) = \arg\max_{\boldsymbol\theta} \ell(\boldsymbol\theta),
\qquad \ell(\boldsymbol\theta) = \ln L(\boldsymbol\theta).
\]

\textbf{چرا لگاریتم؟} 1. با مشاهدات مستقل، درست‌نمایی یک \textbf{حاصل‌ضرب}
است، \(L = \prod_i p(y_i\mid\boldsymbol\theta)\)، که در کامپیوتر حتی
برای چند صد \lr{trial} به 0 سرریز \lr{(underflow)} می‌کند. لگاریتم آن را
به یک \textbf{مجموع} تبدیل می‌کند،
\(\ell = \sum_i \ln p(y_i\mid\boldsymbol\theta)\). 2. لگاریتم اکیداً
صعودی است، پس بیشینه‌کننده‌ی یکسانی دارد. 3. مشتق گرفتن از مجموع راحت‌تر
است.

\(\mathrm{NLL}(\boldsymbol\theta) = -\ell(\boldsymbol\theta)\)

\textbf{مثال تحلیلی (دوجمله‌ای).}
\(\ell(\theta) = \text{\lr{const}} + y\ln\theta + (n-y)\ln(1-\theta)\). با
صفر قرار دادن مشتق:

\[
\frac{d\ell}{d\theta} = \frac{y}{\theta} - \frac{n-y}{1-\theta} = 0 \quad\Rightarrow\quad \hat\theta = \frac{y}{n}.
\]

\lr{MLE} یک احتمال همان نسبت مشاهده‌شده است، که خیال‌مان را راحت می‌کند. در
مدل‌های شناختی، احتمال پیش‌بینی‌شده \emph{تابعی} از پارامترهاست،
\(\theta_i = f(x_i;\boldsymbol\theta)\)، و معمولاً به یک بهینه‌ساز عددی
نیاز داریم.

\subsubsection{3. درست‌نمایی دوجمله‌ای برای منحنی‌های فراموشی}

با \(y_i\) آیتم یادآوری‌شده از \(n_i\) در فاصله‌ی \(t_i\) و فاصله‌های
مستقل:

\[
\ell(\boldsymbol\theta) = \sum_{i=1}^{N}\Big[\ln\binom{n_i}{y_i} + y_i \ln f(t_i;\boldsymbol\theta) + (n_i - y_i)\ln\big(1 - f(t_i;\boldsymbol\theta)\big)\Big].
\]

\begin{itemize}
\tightlist
\item
  مشاهداتی که پیش‌بینی می‌شود نزدیک 0 یا 1 باشند واریانس کمی دارند، پس
  پیش‌بینی‌ای که آن‌ها را از دست بدهد بیشتر جریمه می‌شود. \lr{MLE} همان
  وزن‌دهی واریانسی کمترین مربعات وزنی هفته‌ی 3 را \emph{خودکار} انجام
  می‌دهد.
\item
  \textbf{ایمنی عددی:} اگر مدل جایی که \(y > 0\) است \(f = 0\) پیش‌بینی
  کند، \(\ln 0 = -\infty\) می‌شود. پیش‌بینی‌ها را به
  \([\epsilon, 1-\epsilon]\) محدود \lr{(clip)} کنید، مثلاً با
  \(\epsilon = 10^{-10}\).
\end{itemize}

\textbf{برنولی در برابر دوجمله‌ای.} اگر پاسخ‌های تکی 0/1 را نگه داریم،
\(\ell = \sum_{\text{\lr{trials}}} [y\ln f + (1-y)\ln(1-f)]\). این فقط با
ضرایب دوجمله‌ای ثابت با نسخه‌ی دوجمله‌ای فرق دارد، پس همان \lr{MLE} را
می‌دهد.

\paragraph{خواندن ریاضی}

لگاریتم درست‌نمایی فراموشی همان لگاریتم درست‌نمایی یک نسبت (بخش 2) است که
برای هر فاصله‌ی زمانی یک بار نوشته و جمع شده است:

\begin{itemize}
\tightlist
\item
  برای فاصله‌ی \(i\): جمله‌ی \(y_i \ln f(t_i)\) به پیش‌بینی احتمال یادآوری
  بالا در جایی که آیتم‌های زیادی به یاد آمده‌اند پاداش می‌دهد؛
  \((n_i - y_i)\ln(1 - f(t_i))\) به پیش‌بینی احتمال پایین در جایی که
  آیتم‌های زیادی فراموش شده‌اند پاداش می‌دهد.
\item
  جمع زدن روی فاصله‌ها مجاز است چون فرض می‌کنیم فاصله‌ها \textbf{مستقل}‌اند:
  لگاریتم حاصل‌ضرب احتمال‌ها برابر جمع لگاریتم‌هاست.
\item
  چرا وزن‌دهی خودکار است: اگر مدل \(f = 0.95\) پیش‌بینی کند ولی فقط 40 از
  50 به یاد آمده باشد، جمله‌ی \(10 \ln(0.05) \approx -30\) خیلی پرهزینه
  است. نزدیک \(f = 0.5\) همان مقدار خطا هزینه‌ی خیلی کمتری دارد. به
  نقطه‌هایی که \emph{باید} دقیق باشند بیشتر اعتماد می‌شود.
\end{itemize}

\subsubsection{4. حالت گاوسی: کمترین مربعات همان بیشینه درست‌نمایی است}

فرض کنید \(y_i = f(x_i;\boldsymbol\theta) + \varepsilon_i\) با
\(\varepsilon_i \sim \mathcal{N}(0,\allowbreak \sigma)\)، مستقل و با \(\sigma\)
\textbf{یکسان}. آن‌وقت

\[
\ell(\boldsymbol\theta, \sigma) = \sum_{i=1}^{N}\ln\left[\frac{1}{\sqrt{2\pi}\sigma}\exp\left(-\frac{(y_i - f(x_i;\boldsymbol\theta))^2}{2\sigma^2}\right)\right]
= -\frac{N}{2}\ln(2\pi\sigma^2) - \frac{\mathrm{SSE}(\boldsymbol\theta)}{2\sigma^2}.
\]

برای هر \(\sigma\) ثابت، بیشینه کردن \(\ell\) روی \(\boldsymbol\theta\)
یعنی کمینه کردن \lr{SSE.} پس \textbf{کمترین مربعات = \lr{MLE} تحت نویز
گاوسی مستقل و با واریانس برابر}. بیشینه کردن روی \(\sigma\) هم
\(\hat\sigma^2 = \mathrm{SSE}/N\) را می‌دهد.

این توضیح می‌دهد کمترین مربعات \emph{کی} مناسب است (نویز پیوسته، تقریباً
متقارن، با واریانس ثابت) و کی نیست (شمارش‌ها، نسبت‌ها، انتخاب‌ها، زمان‌های
واکنش چوله).

\begin{quote}
\textbf{پیشرفته: ویژگی‌های \lr{MLE.}} تحت شرایط منظم بودن (مدل
\lr{identifiable}، پارامترهای واقعی نه روی مرز، \(N\) بزرگ)، \lr{MLE}
\emph{سازگار} \lr{(consistent:} به مقدار واقعی همگرا می‌شود)، \emph{به‌طور
مجانبی نرمال} و \emph{به‌طور مجانبی کارا} است (هیچ برآوردگر سازگاری در حد
واریانس کمتری ندارد). واریانس نمونه‌گیری آن با معکوس \textbf{اطلاع فیشر}
\lr{(Fisher information)}، یعنی انحنای \(-\ell\) در کمینه‌اش، تقریب زده
می‌شود. هفته‌ی 5 دقیقاً از همین برای محاسبه‌ی خطاهای استاندارد استفاده
می‌کند. در نمونه‌های کوچک \lr{MLE} ممکن است سوگیری داشته باشد (مثلاً
\(\hat\sigma^2 = \mathrm{SSE}/N\) واریانس \(\sigma^2\) را کمتر برآورد
می‌کند).
\end{quote}

\subsubsection{\texorpdfstring{5. برازش مطلق: انحراف
\(G^2\)}{5. برازش مطلق: انحراف G\^{}2}}

مقدار لگاریتم درست‌نمایی به‌تنهایی معنای مطلقی ندارد. برای قضاوت درباره‌اش،
مدل را با \textbf{مدل اشباع‌شده} \lr{(saturated)} مقایسه کنید که برای هر
خانه‌ی داده یک احتمال آزاد دارد و بنابراین هر نسبت مشاهده‌شده را دقیقاً
بازتولید می‌کند (\(\hat f_i = y_i/n_i\)). \textbf{انحراف} \lr{(deviance)}
برابر است با

\[
G^2 = 2\big[\ell_{\text{\lr{saturated}}} - \ell(\hat{\boldsymbol\theta})\big]
= 2\sum_i \left[y_i \ln\frac{y_i}{n_i \hat f_i} + (n_i - y_i)\ln\frac{n_i - y_i}{n_i(1-\hat f_i)}\right],
\]

با \(0 \ln 0 = 0\). اگر مدل درست باشد و شمارش‌ها خیلی کوچک نباشند،
\(G^2\) تقریباً توزیع \(\chi^2\) با
\(\text{\lr{df}} = (\text{تعداد خانه‌ها}) - k\) دارد. \(G^2\) بزرگ (\(p\)
کوچک) یعنی عدم برازش مدل \emph{بیش از آن است که نویز نمونه‌گیری بتواند
توضیح دهد}.

هشدارها: با \(N\) خیلی بزرگ، هر عدم برازش کوچکی «معنادار» می‌شود؛ با
\(N\) کوچک، عدم برازش واقعی ممکن است دیده نشود. یک \(G^2\) غیرمعنادار
\textbf{ثابت نمی‌کند} که مدل درست است (هفته‌ی 6).

\paragraph{خواندن ریاضی}

\begin{itemize}
\tightlist
\item
  \textbf{مدل اشباع‌شده \lr{(saturated)}} مدل «تقلب‌کننده»‌ای است که هر
  نسبت مشاهده‌شده را دقیقاً پیش‌بینی می‌کند. لگاریتم درست‌نمایی آن بهترین
  مقداری است که هر مدلی ممکن است روی این داده به دست آورد.
\item
  \(G^2 = 2(\ell_{\text{\lr{saturated}}} - \ell(\hat{\boldsymbol\theta}))\)
  اندازه می‌گیرد مدل ما چقدر از این سقف فاصله دارد. \(G^2 = 0\) یعنی
  برازش کامل.
\item
  \textbf{درجه‌ی آزادی \lr{(df)}} = تعداد خانه‌های داده \ensuremath{-} تعداد پارامترهای
  برازش‌شده. با 8 فاصله و 2 پارامتر، \lr{df} = 6.
\item
  اگر مدل درست باشد، \(G^2\) مثل یک متغیر \(\chi^2\) با همان \lr{df}
  رفتار می‌کند که میانگینش \lr{df} است. پس \(G^2 \approx 6\) با \lr{df} =
  6 مقدار عادی عدم برازش است (فقط نویز)؛ \(G^2 = 26\) با \lr{df} = 6
  (\(p = .0002\)) یعنی مدل به‌طور نظام‌مند داده را از دست می‌دهد.
\end{itemize}

\subsubsection{6. مدل‌های مقیدشده به‌عنوان فرضیه}

\lr{MLE} برازش \textbf{چند شرط هم‌زمان با بعضی پارامترهای مشترک} را آسان
می‌کند. مثال: یک آزمایش \lr{SDT} با سه شرط پاداش (پاداش برای پاسخ «بله»،
خنثی، پاداش برای «نه»).

\begin{longtable}[]{@{}
  >{\raggedright\arraybackslash}p{(\linewidth - 6\tabcolsep) * \real{0.2500}}
  >{\raggedright\arraybackslash}p{(\linewidth - 6\tabcolsep) * \real{0.2500}}
  >{\raggedright\arraybackslash}p{(\linewidth - 6\tabcolsep) * \real{0.2500}}
  >{\raggedright\arraybackslash}p{(\linewidth - 6\tabcolsep) * \real{0.2500}}@{}}
\caption{\centering 6. مدل‌های مقیدشده به‌عنوان فرضیه}\tabularnewline
\toprule\noalign{}
\begin{minipage}[b]{\linewidth}\raggedright
مدل
\end{minipage} & \begin{minipage}[b]{\linewidth}\raggedright
پارامترهای آزاد
\end{minipage} & \begin{minipage}[b]{\linewidth}\raggedright
\(k\)
\end{minipage} & \begin{minipage}[b]{\linewidth}\raggedright
فرضیه
\end{minipage} \\
\midrule\noalign{}
\endfirsthead
\toprule\noalign{}
\begin{minipage}[b]{\linewidth}\raggedright
مدل
\end{minipage} & \begin{minipage}[b]{\linewidth}\raggedright
پارامترهای آزاد
\end{minipage} & \begin{minipage}[b]{\linewidth}\raggedright
\(k\)
\end{minipage} & \begin{minipage}[b]{\linewidth}\raggedright
فرضیه
\end{minipage} \\
\midrule\noalign{}
\endhead
\bottomrule\noalign{}
\endlastfoot
حساسیت مشترک & \(d'\)؛ \(c_1, c_2, c_3\) & 4 & پاداش فقط \emph{سوگیری}
را تغییر می‌دهد \\
جداگانه & \(d'_1, d'_2, d'_3\)؛ \(c_1, c_2, c_3\) & 6 & پاداش ممکن است
\emph{حساسیت} را هم تغییر دهد (این مدل اشباع‌شده برای 6 نرخ است) \\
\end{longtable}

لگاریتم درست‌نمایی مشترک، مجموع لگاریتم درست‌نمایی‌های دوجمله‌ای \lr{hit}ها
و \lr{false alarm}ها در همه‌ی شرایط است. مقایسه‌ی لگاریتم درست‌نمایی
بیشینه‌ی دو مدل یک سؤال واقعی را آزمون می‌کند. برای مدل‌های \lr{nested} مثل
این‌ها، \(2(\ell_{6} - \ell_{4})\) تقریباً \(\chi^2\) با 2 درجه‌ی آزادی
است: \textbf{آزمون نسبت درست‌نمایی} که در هفته‌ی 7 کامل می‌شود.

\begin{quote}
\textbf{پیشرفته: \lr{MLE} در \lr{SDT} و فرمول‌ها.} برای یک شرط تنها،
\lr{MLE} \((d', c)\) از شمارش‌ها دقیقاً همان چیزی است که فرمول‌های هفته‌ی 2
\emph{بدون} تصحیح می‌دهند، چون مدل اشباع‌شده است (2 پارامتر، 2 نرخ) و
\lr{MLE} نرخ‌های مشاهده‌شده را بازتولید می‌کند. تصحیح \lr{log-linear}،
\lr{MLE} \emph{نیست}؛ بیشتر به یک برآورد بیزی با پیشین ضعیف شبیه است
(هفته‌ی 9).
\end{quote}

\paragraph{برای یک مهندس}

\begin{itemize}
\tightlist
\item
  این دقیقاً همان \textbf{تخمین بیشینه‌ی درست‌نمایی} در نظریه‌ی تخمین است
  (تخمین \lr{ML} کانال یا فرکانس، آشکارسازی دنباله با \lr{ML).} تنها
  تفاوت مدل نویز است: شمارش دوجمله‌ای به جای \lr{AWGN.}
\item
  \textbf{\lr{LS} = \lr{ML} تحت \lr{AWGN}} (بخش 4) نتیجه‌ای است که قبلاً
  از آن استفاده کرده‌اید: با نویز سفید گاوسی با واریانس معلوم و برابر،
  تخمین‌گر \lr{ML} خطای مربعی را کمینه می‌کند.
\item
  \textbf{اطلاعات فیشر \lr{(Fisher information)}} و واریانس مجانبی
  \lr{MLE} (کادر پیشرفته) همان کمیت‌های \textbf{کران پایین
  \lr{Cramér\lr{--}Rao}} هستند. هفته‌ی 5 از آن‌ها برای محاسبه‌ی خطای استاندارد
  استفاده می‌کند.
\item
  مقایسه‌ی یک مدل مقید و یک مدل آزاد با
  \(2(\ell_{\text{\lr{free}}} - \ell_{\text{\lr{constrained}}})\) همان
  \textbf{آزمون نسبت درست‌نمایی تعمیم‌یافته \lr{(GLRT)}} است: هر دو فرضیه
  پارامترهای مجهول دارند که با تخمین \lr{ML} جایگزین می‌شوند.
\item
  \textbf{پاداش‌ها و جابه‌جایی ملاک:} در آشکارسازی بیزی، آستانه‌ی بهینه با
  هزینه‌ها و احتمال‌های پیشین جابه‌جا می‌شود. بخش 6 می‌آزماید که آیا
  \emph{آدم‌ها} وقتی پاداش‌ها عوض می‌شوند فقط آستانه‌شان (\(c\)) را جابه‌جا
  می‌کنند و \lr{«SNR»}شان (\(d'\)) ثابت می‌ماند.
\end{itemize}

\paragraph{برای مدرس}

\begin{itemize}
\tightlist
\item
  \textbf{جمله‌ای که باید تکرار کنید:} «درست‌نمایی احتمال داده‌ای است که
  مشاهده کرده‌ایم، به‌عنوان تابعی از پارامترها». آن را روی تخته بنویسید و
  تا آخر جلسه نگه دارید.
\item
  \textbf{رسمش کنید.} \(L(\theta)\) را برای 7 از 10 (مثال حل‌شده، بخش 1)
  به شکل منحنی روی \(\theta\) از 0 تا 1 بکشید. بپرسید: قله کجاست؟ آیا
  مساحت زیر منحنی 1 است؟ (نه.) بعد آن را برای 70 از 100 بکشید: همان قله،
  منحنی خیلی باریک‌تر. این باریکی همان دقت \lr{(precision)} است، موضوع
  هفته‌ی 5.
\item
  \textbf{از پیش‌زمینه‌ی خودتان استفاده کنید.} خیلی از دانشجوها \lr{MLE}
  را انتزاعی می‌دانند؛ شما می‌توانید آن را به‌عنوان «روش تخمینی که مهندس‌ها
  وقتی نویز گاوسی نیست به کار می‌برند» معرفی کنید.
\item
  \textbf{سؤال‌های رایج.}

  \begin{itemize}
  \tightlist
  \item
    \emph{«آیا درست‌نمایی احتمال درست بودن \(\theta\) است؟»} نه. برای آن
    به \lr{prior} نیاز داریم (هفته‌ی 9).
  \item
    \emph{«چرا \lr{NLL} را کمینه می‌کنیم به جای اینکه \(L\) را بیشینه
    کنیم؟»} جواب یکی است، ولی جمع لگاریتم‌ها از \lr{underflow} عددی
    جلوگیری می‌کند و بهینه‌سازها برای کمینه کردن نوشته شده‌اند.
  \item
    \emph{«آیا \(G^2\) غیرمعنادار ثابت می‌کند مدل درست است؟»} نه؛ فقط
    یعنی عدم برازش بیشتر از نویز نیست، و با داده‌ی کم ممکن است عدم برازش
    واقعی پنهان بماند.
  \end{itemize}
\item
  \textbf{مدیریت زمان.} بخش 4 (حالت گاوسی) را می‌شود سریع به‌عنوان
  \lr{«LS} حالت خاصی است» نشان داد. وقت صرفه‌جویی‌شده را روی بخش 5
  (\(G^2\)) بگذارید.
\end{itemize}

\begin{center}\rule{0.5\linewidth}{0.5pt}\end{center}

\Needspace{11\baselineskip}

\subsection{مثال حل‌شده}

\subsubsection{بخش 1: درست‌نمایی یک نسبت تنها}

\(y = 7\) یادآوری از \(n = 10\):

\begin{longtable}[]{@{}
  >{\raggedright\arraybackslash}p{(\linewidth - 8\tabcolsep) * \real{0.2000}}
  >{\raggedright\arraybackslash}p{(\linewidth - 8\tabcolsep) * \real{0.2000}}
  >{\raggedright\arraybackslash}p{(\linewidth - 8\tabcolsep) * \real{0.2000}}
  >{\raggedright\arraybackslash}p{(\linewidth - 8\tabcolsep) * \real{0.2000}}
  >{\raggedright\arraybackslash}p{(\linewidth - 8\tabcolsep) * \real{0.2000}}@{}}
\caption{\centering بخش 1: درست‌نمایی یک نسبت تنها}\tabularnewline
\toprule\noalign{}
\begin{minipage}[b]{\linewidth}\raggedright
\(\theta\)
\end{minipage} & \begin{minipage}[b]{\linewidth}\raggedright
0.5
\end{minipage} & \begin{minipage}[b]{\linewidth}\raggedright
0.6
\end{minipage} & \begin{minipage}[b]{\linewidth}\raggedright
\textbf{0.7}
\end{minipage} & \begin{minipage}[b]{\linewidth}\raggedright
0.8
\end{minipage} \\
\midrule\noalign{}
\endfirsthead
\toprule\noalign{}
\begin{minipage}[b]{\linewidth}\raggedright
\(\theta\)
\end{minipage} & \begin{minipage}[b]{\linewidth}\raggedright
0.5
\end{minipage} & \begin{minipage}[b]{\linewidth}\raggedright
0.6
\end{minipage} & \begin{minipage}[b]{\linewidth}\raggedright
\textbf{0.7}
\end{minipage} & \begin{minipage}[b]{\linewidth}\raggedright
0.8
\end{minipage} \\
\midrule\noalign{}
\endhead
\bottomrule\noalign{}
\endlastfoot
\(L(\theta) = \binom{10}{7}\theta^7(1-\theta)^3\) & 0.117 & 0.215 &
\textbf{0.267} & 0.201 \\
\(\ell(\theta)\) & \ensuremath{-}2.144 & \ensuremath{-}1.537 & \textbf{\ensuremath{-}1.321} & \ensuremath{-}1.603 \\
\end{longtable}

بیشینه در \(\hat\theta = 7/10 = 0.7\) است. دقت کنید که مقادیر درست‌نمایی
روی \(\theta\) جمعشان 1 نمی‌شود: درست‌نمایی توزیع احتمال روی پارامترها
نیست.

\subsubsection{بخش 2: یک آزمایش فراموشی دوفاصله‌ای}

\(y = (30, 13)\) آیتم یادآوری‌شده از \(n = 50\) در \(t = (1, 7)\) روز.
مدل: \(f(t) = A e^{-\lambda t}\).

\textbf{قدم 1: \lr{NLL} برای یک نامزد.} برای
\((A, \lambda) = (0.70, 0.15)\): \(f(1) = 0.7e^{-0.15} = 0.6025\)،
\(f(7) = 0.7e^{-1.05} = 0.2450\).

\[
\ell = \ln\mathrm{Binom}(30; 50, 0.6025) + \ln\mathrm{Binom}(13; 50, 0.2450) = -2.167 + (-2.088) = -4.255,
\]

پس \lr{NLL} \(= 4.255\).

\textbf{قدم 2: مقایسه.} برای \((0.70, 0.12)\): \(f = (0.621, 0.302)\) و
\lr{NLL} \(= 4.488\) که بدتر است. نامزد اول داده‌ها را محتمل‌تر می‌کند.

\textbf{قدم 3: \lr{MLE.}} با 2 پارامتر و 2 داده، مدل می‌تواند هر دو نسبت
مشاهده‌شده را دقیقاً بازتولید کند: \(A e^{-\lambda} = 0.60\) و
\(A e^{-7\lambda} = 0.26\). تقسیم این دو می‌دهد
\(e^{6\lambda} = 0.60/0.26\)، پس \(\hat\lambda = \ln(2.308)/6 = 0.139\)
و \(\hat A = 0.60\,e^{0.139} = 0.690\)، با \lr{NLL} \(= 4.225\) (کمینه).
یک بهینه‌ساز عددی همین مقادیر را می‌دهد.

\textbf{تفسیر.} \lr{MLE} می‌گوید این شرکت‌کننده در شروع حدود 69\% را به
یاد می‌آورد و هر روز حدود 13\% از باقی‌مانده را از دست می‌دهد
(\(1 - e^{-0.139}\)). چون اینجا \(k = N\) است، \(G^2 = 0\): مدل
\emph{نمی‌تواند} بد برازش شود، پس این برازش کامل هیچ شاهدی برای شکل نمایی
نیست. فقط داده‌هایی با تعداد خانه‌ی بیشتر از تعداد پارامترها می‌توانند یک
مدل را آزمون کنند.

\begin{center}\rule{0.5\linewidth}{0.5pt}\end{center}

\Needspace{11\baselineskip}

\subsection{\texorpdfstring{\lr{lab} (2 ساعت)}{lab (2 ساعت)}}

\subsubsection{اهداف}

\begin{enumerate}
\def\labelenumi{\arabic{enumi}.}
\tightlist
\item
  رسم یک تابع درست‌نمایی و دیدن اینکه چرا با لگاریتم درست‌نمایی کار
  می‌کنیم.
\item
  برازش مدل‌های فراموشی نمایی و توانی با \lr{MLE}؛ مقایسه با کمترین
  مربعات؛ ارزیابی با \(G^2\).
\item
  وارسی عددی اینکه \lr{MLE} گاوسی با کمترین مربعات برابر است.
\item
  برازش یک مدل \lr{SDT} چندشرطی مقیدشده و \textbf{تفسیر} آنچه درباره‌ی
  پاداش‌ها می‌گوید.
\end{enumerate}

\textbf{زمان‌بندی:} بخش \lr{A} 15 دقیقه · بخش \lr{B} 35 دقیقه · بخش
\lr{C} 15 دقیقه · بخش \lr{D} 35 دقیقه · بخش \lr{E} 20 دقیقه (تمرین‌ها).

\subsubsection{\texorpdfstring{بخش \lr{B: MLE} برای منحنی‌های فراموشی (35
دقیقه)}{بخش B: MLE برای منحنی‌های فراموشی (35 دقیقه)}}

\textbf{به این‌ها دقت کنید:} - هر روش در معیار \emph{خودش} برنده است:
\lr{MLE} کمترین \lr{NLL} و \lr{LS} کمترین \lr{SSE} را دارد. - برای مدل
توانی (مدل واقعی) دو روش خیلی به هم نزدیک‌اند. - برای مدل نمایی دو روش
بیشتر با هم فرق دارند (نرخ زوال 0.046 در برابر 0.057). درست‌نمایی
دوجمله‌ای به نقاطی که نزدیک کف دقیق اندازه‌گیری شده‌اند (فاصله‌های طولانی)
وزن بیشتری می‌دهد، پس منحنی \lr{ML} کندتر زوال پیدا می‌کند تا به آن‌ها
نزدیک‌تر بماند.

\begin{quote}
خیلی بدتر از کمینه‌ی واقعی 29.8 (حتی بدتر از \lr{NLL} برآورد \lr{LS).}
روش‌های سیمپلکس کران‌دار وقتی رأس‌ها روی یک کران فشرده شوند ممکن است
تباهیده \lr{(degenerate)} شوند. دو بررسی سریع این را لو می‌دهند: برآوردی
که دقیقاً روی یک کران نشسته، و یک \lr{«MLE»} که \lr{NLL} آن از مقادیر
پارامتر دیگری که امتحان کرده‌اید \emph{بیشتر} است. \lr{multi-start}
(هفته‌ی 3) درمان کلی آن است.
\end{quote}

حالا برازش مطلق را با انحراف ارزیابی کنید:

\textbf{تفسیر:} عدم برازش مدل نمایی خیلی بیشتر از آن است که نویز
دوجمله‌ای اجازه می‌دهد (\(p = .0002\)): می‌توانیم آن را به‌عنوان توصیف این
داده‌ها \textbf{رد کنیم}. عدم برازش مدل توانی با نویز نمونه‌گیری سازگار
است. این ثابت نمی‌کند که مدل درستی است؛ یعنی این داده‌ها دلیلی برای رد آن
نمی‌دهند.

\subsubsection{\texorpdfstring{بخش \lr{C: MLE} گاوسی همان کمترین مربعات
است (15 دقیقه)}{بخش C: MLE گاوسی همان کمترین مربعات است (15 دقیقه)}}

نسبت‌ها را پیوسته در نظر بگیرید و نویز گاوسی با \(\sigma\) آزاد (پارامتر
سوم) فرض کنید. برآوردهای \(A\) و \(d\) باید با کمترین مربعات یکی باشند و
\(\hat\sigma\) باید برابر \(\sqrt{\mathrm{SSE}/N}\) باشد.

\subsubsection{\texorpdfstring{بخش \lr{D:} یک مدل \lr{SDT} چندشرطی
مقیدشده (35 دقیقه)}{بخش D: یک مدل SDT چندشرطی مقیدشده (35 دقیقه)}}

\textbf{آزمایش.} شرکت‌کننده‌ها یک صدای ضعیف را تحت سه طرح پاداش تشخیص
می‌دهند: \textbf{\lr{liberal}} (پاداش برای \lr{hit)}، \textbf{خنثی} و
\textbf{محافظه‌کار} (جریمه برای \lr{false alarm).} در هر شرط 80
\lr{trial} علامت و 80 \lr{trial} نویز. یک شرکت‌کننده را شبیه‌سازی می‌کنیم
که حساسیتش در همه‌ی شرایط یکسان است و ملاکش با پاداش جابه‌جا می‌شود.

با 80 \lr{trial} در هر خانه این یک خطای نمونه‌گیری عادی است. هفته‌ی 5 نشان
می‌دهد چطور دور آن بازه‌ی اطمینان بسازیم.

آیا مقید کردن \(d'\) به برابری در همه‌ی شرایط پاداش درست‌نمایی زیادی را
هزینه می‌کند؟ درباره‌ی اثر پاداش‌ها چه نتیجه‌ای می‌گیرید؟ 2. برآوردهای \(d'\)
جداگانه چند دهم با هم فرق دارند. چرا این شاهدی بر اینکه پاداش حساسیت را
تغییر می‌دهد نیست؟ 3. بررسی کنید که مدل \(d'\) جداگانه نرخ‌های \lr{hit} و
\lr{false alarm} مشاهده‌شده را دقیقاً بازتولید می‌کند. چرا باید این‌طور
باشد؟

\subsubsection{تمرین‌ها (به ترتیب سختی)}

\(y = 13\) · \(n = 40\)

\textbf{تمرین 2 \ensuremath{\star}.} در بخش \lr{B}، هر دو مدل را با درست‌نمایی
\textbf{برنولی} روی داده‌ی سطح آیتم دوباره برازش کنید (هر شمارش را به
پاسخ‌های 0/1 باز کنید). نشان دهید برآوردها با برازش‌های دوجمله‌ای یکسان‌اند
و \lr{NLL}ها به اندازه‌ی یک ثابت فرق دارند. آن ثابت چیست؟

هر دو مدل را با \lr{MLE} برازش کنید و درباره‌ی آنچه بر سر برآوردها و
\(G^2\) می‌آید توضیح دهید.

\textbf{تمرین 4 \ensuremath{\star}\ensuremath{\star} (یک مدل \lr{SDT} سوم).} مدلی را برازش کنید که در آن
ملاک در همه‌ی شرایط \emph{یکسان} است ولی \(d'\) فرق دارد (\(k = 4\)).
\lr{NLL} آن را با مدل \(d'\) مشترک مقایسه کنید. هر دو \(k = 4\) دارند؛
چرا اینجا می‌توانید مستقیماً با \lr{NLL} مقایسه کنید؟

\(A = 0.9\) · \(d = 0.5\) · \(n = 50\) با \(n = 10\) تکرار کنید. (آمادگی
خوبی برای \lr{A2} است.)

\subsubsection{\texorpdfstring{\lr{lab check-off}}{lab check-off}}

نشان دهید: (1) جدول \lr{LS} در برابر \lr{MLE} و نتایج \(G^2\) از بخش
\lr{B}، (2) مقایسه‌ی گاوسی در برابر \lr{LS} از بخش \lr{C}، و (3)
پاسخ‌هایتان به سؤال‌های تفسیری بخش \lr{D.}

\begin{center}\rule{0.5\linewidth}{0.5pt}\end{center}

\subsection{\texorpdfstring{ارتباط با هوش مصنوعی و
\lr{LLM}ها}{ارتباط با هوش مصنوعی و LLMها}}

\textbf{\lr{LLM}ها با بیشینه‌ی درست‌نمایی آموزش می‌بینند.}
\lr{«cross-entropy loss»} که برای آموزش همه‌ی مدل‌های زبانی استفاده می‌شود،
همان منفی لگاریتم درست‌نمایی \lr{token} بعدی است که روی یک \lr{corpus}
متنی عظیم میانگین گرفته شده. پس آموزش یک \lr{LLM} یک مسئله‌ی \lr{MLE} با
میلیاردها پارامتر است که به جای \lr{Nelder\lr{--}Mead} با \lr{gradient
descent} حل می‌شود. دو کمیتی که در مقاله‌های هوش مصنوعی می‌بینید همان
کمیت‌های این هفته در لباس دیگرند: \textbf{\lr{loss}} = \lr{NLL} برای هر
\lr{token}، و \textbf{\lr{perplexity}} = \(e^{\text{\lr{NLL per token}}}\)،
یعنی تعداد مؤثر گزینه‌های هم‌احتمالی که مدل بینشان «سردرگم» است.

ساده‌ترین مدل زبانی، یعنی \textbf{مدل \lr{bigram}}، یک \lr{MLE} بسته
دارد: \(P(w_2\mid w_1)\) = \lr{count(}\(w_1 w_2\)\lr{)/count(}\(w_1\)).
این مدل یک مشکل آشنا را نشان می‌دهد. هر جفت کلمه‌ای که در آموزش دیده نشده
احتمال 0 می‌گیرد، پس فقط یک جفت تازه درست‌نمایی کل جمله را 0 می‌کند
\lr{(NLL} = \ensuremath{\infty}). هموارسازی \lr{(smoothing)} این را حل می‌کند، همان‌طور که
تصحیح \lr{log-linear} هفته‌ی 2 مشکل نرخ \lr{hit} 1 را حل کرد.

با دست بررسی کنید که \lr{MLE} برای \(P(\text{\lr{cat}}\mid\text{\lr{the}})\) برابر
1/3 است. - \ensuremath{\star}\ensuremath{\star} هموارسازی جمله‌ی \emph{آموزشی} را کم‌احتمال‌تر می‌کند
\lr{(NLL} از 3.6 به 9.5) ولی جمله‌ی تازه را ممکن می‌کند. \lr{add-one} را
با \lr{add-}\(k\) عوض کنید و \(k\) را با \lr{cross-validation
leave-one-sentence-out} انتخاب کنید (پیش‌نمایشی از هفته‌ی 6). - \ensuremath{\star}\ensuremath{\star} بحث:
\lr{MLE} یک مدل \lr{bigram corpus} آموزشی را تا جای ممکن بازتولید می‌کند.
چرا «بیشینه کردن درست‌نمایی روی متن اینترنت» با «یاد گرفتن راست‌گویی» یکی
نیست؟

\begin{center}\rule{0.5\linewidth}{0.5pt}\end{center}

\subsection{اشتباهات رایج و برداشت‌های نادرست}

\begin{enumerate}
\def\labelenumi{\arabic{enumi}.}
\tightlist
\item
  \textbf{«درست‌نمایی احتمال درست بودن پارامترهاست.»} احتمال (چگالی)
  \emph{داده} به شرط پارامترهاست. انتگرالش روی \(\boldsymbol\theta\)
  برابر 1 نیست.
\item
  \textbf{ضرب کردن احتمال‌ها به‌جای جمع کردن لگاریتم‌ها.} حاصل‌ضرب برای
  تعداد \lr{trial} واقعی به 0 سرریز می‌کند (بخش \lr{A).}
\item
  فراموش کردن علامت منفی بهینه‌ساز را به \emph{بدترین} پارامترها می‌فرستد.
\item
  \textbf{محدود نکردن احتمال‌ها.} احتمال‌های پیش‌بینی‌شده‌ی دقیقاً 0 یا 1
  لگاریتم درست‌نمایی \(-\infty\) می‌دهند و بهینه‌ساز را از کار می‌اندازند.
\item
  \textbf{استفاده‌ی پیش‌فرض از کمترین مربعات برای انتخاب‌ها یا شمارش‌ها.}
  ساختار نویز داده را نادیده می‌گیرد. از درست‌نمایی‌ای استفاده کنید که با
  نحوه‌ی تولید داده جور باشد.
\item
  \textbf{خواندن \(G^2\) غیرمعنادار به‌عنوان اثبات درستی مدل.} فقط می‌گوید
  عدم برازش با این حجم نمونه به‌طور قابل تشخیصی از نویز نمونه‌گیری بیشتر
  نیست.
\item
  \textbf{فراموش کردن اینکه مدل اشباع‌شده یا دقیقاً تعیین‌شده همیشه کامل
  برازش می‌شود} (مثال حل‌شده، مدل \(d'\) جداگانه در بخش \lr{D).} در این
  حالت برازش کامل بی‌اطلاع است.
\end{enumerate}

\begin{center}\rule{0.5\linewidth}{0.5pt}\end{center}

\subsection{خودآزمایی}

\textbf{پاسخ} احتمال: پارامترها ثابت، داده متغیر، و جمع/انتگرالش روی
داده 1 است. درست‌نمایی: داده روی مقادیر مشاهده‌شده ثابت، پارامترها متغیر،
و لازم نیست انتگرالش روی پارامترها 1 باشد.

\textbf{2.} دو دلیل برای کار با لگاریتم درست‌نمایی بیاورید. \textbf{پاسخ}
(1) عددی: حاصل‌ضرب احتمال‌های زیاد سرریز می‌کند ولی مجموع لگاریتم‌ها نه. (2)
لگاریتم \lr{monotone} است، پس بیشینه‌کننده عوض نمی‌شود، و مشتق‌گیری و
بهینه‌سازی مجموع‌ها راحت‌تر است.

\textbf{3.} تحت چه فرض‌هایی کمترین مربعات با بیشینه درست‌نمایی معادل است؟
گام کلیدی را نشان دهید. \textbf{پاسخ} خطاهای گاوسی مستقل با واریانس
برابر \(\sigma^2\). آن‌وقت
\(\ell = -\tfrac{N}{2}\ln(2\pi\sigma^2) - \mathrm{SSE}/(2\sigma^2)\)، پس
برای هر \(\sigma\)، بیشینه کردن \(\ell\) روی \(\boldsymbol\theta\) همان
کمینه کردن \lr{SSE} است.

\textbf{4.} یک مدل 2 پارامتری روی 8 خانه‌ی دوجمله‌ای برازش شده و
\(G^2 = 4.1\) می‌دهد. درجه‌ی آزادی چند است و چه نتیجه‌ای می‌گیرید؟
\textbf{پاسخ} \lr{df} = 8 \ensuremath{-} 2 = 6. میانگین \(\chi^2_6\) برابر 6 است، پس
4.1 کاملاً در دامنه‌ی مورد انتظار از نویز نمونه‌گیری است
(\(p \approx 0.66\)). شاهدی بر عدم برازش نیست، ولی این درستی مدل را ثابت
نمی‌کند.

\textbf{5.} چرا \lr{MLE} با درست‌نمایی دوجمله‌ای نقاط داده را متفاوت از
کمترین مربعات ساده وزن می‌دهد؟ \textbf{پاسخ} لگاریتم درست‌نمایی دوجمله‌ای
واریانس \(f(1-f)/n\) را منعکس می‌کند: نقاطی که پیش‌بینی می‌شود نزدیک 0 یا 1
باشند (واریانس کم) دقیق‌تر برازش می‌شوند، و نقاط با \(n\) بزرگ‌تر بیشتر
حساب می‌شوند. کمترین مربعات ساده همه‌ی مربع خطاها را یکسان وزن می‌دهد.

\textbf{6.} در آزمایش پاداش، مدل \(d'\) مشترک چه فرضیه‌ی روان‌شناختی را
بیان می‌کند و چطور آزمونش می‌کنید؟ \textbf{پاسخ} اینکه پاداش‌ها فقط سوگیری
پاسخ (ملاک) را تغییر می‌دهند، نه حساسیت ادراکی را. با مقایسه‌ی لگاریتم
درست‌نمایی بیشینه با مدل \(d'\) جداگانه آزمونش کنید:
\(2(\ell_6 - \ell_4) \sim \chi^2_2\) تقریبی (آزمون نسبت درست‌نمایی)، یا
با \lr{AIC/BIC} (هفته‌ی 7).

\textbf{7.} مدلی را برازش می‌کنید و بهینه‌ساز \lr{NLL} = \ensuremath{-}35.2 (منفی)
برمی‌گرداند. آیا اشکالی هست؟ \textbf{پاسخ} نه لزوماً. برای داده‌ی پیوسته،
چگالی‌ها می‌توانند از 1 بیشتر باشند (مثلاً یک گاوسی باریک)، پس لگاریتم
چگالی می‌تواند مثبت و \lr{NLL} منفی باشد. برای داده‌ی گسسته (دوجمله‌ای)
احتمال‌ها \ensuremath{\leq} 1 هستند، پس \lr{NLL} باید \ensuremath{\geq} 0 باشد و مقدار منفی نشانه‌ی باگ
است.

\begin{center}\rule{0.5\linewidth}{0.5pt}\end{center}

\Needspace{11\baselineskip}

\subsection{تکلیف}

\subsubsection{\texorpdfstring{\lr{A2:} کمترین مربعات در برابر بیشینه
درست‌نمایی (1.0 نمره از 20؛ انتشار هفته‌ی 4، موعد ابتدای درس هفته‌ی
6)}{A2: کمترین مربعات در برابر بیشینه درست‌نمایی (1.0 نمره از 20؛ انتشار هفته‌ی 4، موعد ابتدای درس هفته‌ی 6)}}

\textbf{کارها.}

\begin{enumerate}
\def\labelenumi{\arabic{enumi}.}
\tightlist
\item
  \textbf{برازش (2 امتیاز).} مدل‌های نمایی و توانی را روی این داده‌ها با
  (الف) کمترین مربعات روی نسبت‌ها و (ب) \lr{MLE} دوجمله‌ای برازش کنید.
  همه‌ی برآوردها را با \lr{SSE} و \lr{NLL} در یک جدول گزارش کنید.
\item
  \textbf{ارزیابی (2 امتیاز).} \(G^2\)، درجه‌ی آزادی و \(p\) را برای
  برازش \lr{MLE} هر مدل حساب کنید. داده‌ها را با بازه‌های 95\% دوجمله‌ای،
  هر دو منحنی برازش‌شده، و باقی‌مانده‌ها رسم کنید.
\item
  \textbf{مطالعه‌ی شبیه‌سازی (3 امتیاز).} مدل توانی با \(A = 0.95\)،
  \(d = 0.7\) را مدل واقعی در نظر بگیرید. هیستوگرام‌ها را نشان دهید. با
  \(n = 100\) آیتم تکرار کنید و بگویید چه تغییر می‌کند.
\item
  جوابتان را به ساختار نویز داده‌های دوجمله‌ای و به شمارش‌های نزدیک به کف
  (1 تا 2 از 20) در طولانی‌ترین فاصله‌ها ربط دهید. اگر یکی از آن شمارش‌ها 0
  بود، چه بر سر روش خطی‌سازی هفته‌ی 2 می‌آمد؟
\end{enumerate}

برای تحویل با تأخیر، روزی 10\% کسر می‌شود (حداکثر 3 روز).

\begin{center}\rule{0.5\linewidth}{0.5pt}\end{center}

\subsection{مطالعه‌ی بیشتر}

\begin{itemize}
\tightlist
\item
  \textbf{\lr{Farrell} \lr{\&} \lr{Lewandowsky} (2018)، فصل 4.} برآورد
  پارامتر با بیشینه درست‌نمایی: توابع درست‌نمایی برای انواع رایج داده،
  لگاریتم درست‌نمایی و برازش در عمل. \emph{(منبع اصلی.)}
\item
  \textbf{عظیمی‌فر و احمدی (1396)}، \emph{شناسایی آماری الگو}، جلد 1: بخش
  برآورد پارامتر با بیشینه درست‌نمایی، با نگاه یادگیری آماری {[}شماره‌ی
  فصل بررسی شود{]}.
\item
  \textbf{\lr{Lee} \lr{\&} \lr{Wagenmakers} (2014)، فصل 3.} استنباط با توزیع
  دوجمله‌ای. الان برای چیدمان دوجمله‌ای بخوانید؛ در هفته‌ی 9 از سمت بیزی
  دوباره به آن برمی‌گردیم.
\end{itemize}

\emph{پیش‌نمایش هفته‌ی 5:} برآورد بدون عدم قطعیت نصف یک نتیجه است. خطاهای
استاندارد و بازه‌های اطمینان را از انحنای درست‌نمایی و با \lr{bootstrap}
حساب می‌کنیم، و کشف می‌کنیم که چطور میانگین‌گیری از شرکت‌کننده‌ها می‌تواند
«قانونی» بسازد که هیچ فردی از آن پیروی نمی‌کند.
